\documentclass[../../main.tex]{subfiles}

\begin{document}

\section{Geometry in 3-Space}

In this section of the note, we will discuss more geometric
analysis in 3-space with vectors. I personally think it was
very interesting to see different interpretations of vectors
in 3-space because they could be very abstract, yet a fundamental
component that helps analyzing geometry in 3-dimensions.

\subsection{Lines in 3-Space}

Continuing with our prior discussion on vectors, let's discuss
more on lines and planes in 3-space and how they relate to vectors.

First, the components of lines in the coordinate plane can be thought
of as a point and a direction, which would be slope. Given a point
that a line passes through and the slope, or direction, of the line,
we can find the unique line. Similarly, we could do the same for
3-space. The direction in this case would be a nonzero vector.
In other words, given a point and a nonzero vector parallel to
the line, we can find the unique line.

\begin{figure}[H]
\centering
\tdplotsetmaincoords{70}{110}
\begin{tikzpicture}[
        tdplot_main_coords,
        axis/.style={->,thick},
        scale=0.4
    ]
    \draw[axis] (-8,0,0) -- (8,0,0) node[left] {$x$};
    \draw[axis] (0,-7,0) -- (0,8,0) node[right] {$y$};
    \draw[axis] (0,0,-5) -- (0,0,8) node[left] {$z$};

    \draw[axis] (0,0,0) -- (-0.5,3,1.5) node[right]
        {$\langle a, b, c \rangle$};
    \draw[Red] (2,-3,3) -- (1,3,6) node[below] {$l$};
    \node[circle,fill=black,inner sep=0pt,minimum size=1mm]
        at (0,0,0) {};
    \node[circle,fill=black,inner sep=0pt,minimum size=1mm]
        at (1.5, 0, 4.5) {};
    \node[above left] at (1.5, 0, 4.5) {$(x_0, y_0, z_0)$};
\end{tikzpicture}
\end{figure}

Looking at the graph above, we can derive a very important
concept of representing lines in parametric equations.

\begin{important}
A line $l$ in 3-space that passes through the point $(x_0, y_0, z_0)$
and parallel to vector $\mathbf{v} = \langle a, b, c \rangle$ can be
represented as the following parametric equation for
$t \in (-\infty, \infty)$.
\[
x = x_0 + at, \quad
y = y_0 + bt, \quad
z = z_0 + ct
\]
\end{important}

To prove this, we can see that the line actually contains infinitely
many vectors that passes through the point $(x_0, y_0, z_0)$ and
$(x, y, z)$ while being parallel to $\langle a, b, c \rangle$.
In other words, $\langle x - x_0, y - y_0, z - z_0 \rangle =
t\mathbf{v} = \langle ta, tb, tc \rangle$ for some $t$. Comparing
the vectors element-wise, we obtain the parametric equation.

\begin{important}
Another standard way of writing the parametric equation
$x = x_0 + at, y = y_0 + bt, z = z_0 + ct$ is the following.
\[
\langle x, y, z \rangle
= \langle x_0, y_0, z_0 \rangle + t \langle a, b, c \rangle
\]
This is known as a \textbf{vector equation}.
\end{important}

Let's take a look at a quick example.

\begin{problem}{}{}
Find a parametric equation and a vector equation of the line
that passes through the point $(1, 2, 3)$ and parallel to
$\langle 1, 2, 3 \rangle$.
\end{problem}
\begin{solution}
Because the line passes through the point $(1, 2, 3)$ and is
parallel to $\langle 1, 2, 3 \rangle$, a vector equation
can be written as $\langle x, y, z \rangle = \langle 1, 2, 3 \rangle
+ t \langle 1, 2, 3 \rangle$. Similarly, a parametric equation
can be written as $x = 1 + t$, $y = 2 + 2t$, and $z = 3 + 3t$.
\end{solution}

However, this is not the only answer. Because the line is
also parallel to $\langle 1, 2, 3 \rangle$, it is also parallel
to $\langle 2, 4, 6 \rangle$. Using this instead, we get a
vector equation $\langle 1, 2, 3 \rangle + t \langle 2, 4, 6 \rangle$.
This shows that there are different parametrizations and vector
equations for a given line.

One way vector equation can be very helpful than parametric equation
is that we can directly see if two lines are parallel to each other.
If the two vectors that are being scaled by $t$ are scale multiples
of each other, then the two lines are similar. This observation
also leads us to define as new term.

\begin{definition}{}{}
The \textbf{skew} lines are two lines in 3-space that are not
parallel to each other and that do not intersect.
\end{definition}

Unlike 2-space, we can see that not being parallel does not
necessarily imply that the two lines intersect.

\begin{important}
Segments are very similar to lines. When we have to obtain a segment,
then we can just limit our parameter $t$ from the line obtained
from the methods above.
\end{important}

Now that we have took a look at the lines in 3-space, let's discuss
its extension: planes.

\subsection{Planes in 3-Space}

First, let's define the terms \textbf{perpendicular} and
\textbf{parallel} in the context of planes.

\begin{definition}{}{}
A plane $A$ and a vector $\mathbf{v}$ are \textbf{perpendicular}
if any vector lying on the plane $A$ is orthogonal to $\mathbf{v}$.
Moreover, if there exists a vector that is perpendicular to two
planes, then the planes are \textbf{parallel}. Finally, a vector
and a plane is \textbf{parallel} if some normal vector of the plane
is orthogonal to the vector.
\end{definition}

Similar to how we found vector equations of a line, we could
extend our understandings on plane.

\begin{important}{}{}
The plane in 3-space that is perpendicular to a nonzero vector
$\mathbf{v} = \langle a, b, c \rangle$ and contains the point
$(x_0, y_0, z_0)$ can be represented as the following equation.
\[
a(x - x_0) + b(y - y_0) + c(z - z_0) = 0
\]
Such vectors $\mathbf{v}$ that are perpendicular to the plane
are known as \textbf{normal vectors}. Moreover, the equation
above is called the \textbf{point-normal form}.
\end{important}

Similar to how we defined a line, we can notice that the plane
contains all vectors $\langle x - x_0, y - y_0, z - z_0 \rangle$
that is perpendicular to $\mathbf{v}$. Therefore using the dot
product, we know that the vectors are orthogonal if and only
if their dot product is zero.
\[
\langle x - x_0, y - y_0, z - z_0 \rangle \cdot
    \langle a, b, c \rangle
= 0
\]
Expanding the equation, we obtain the equations of a plane.
Besides this point-normal form of the equation of the plane,
we obtain the following form by expanding.
\[
ax + by + cz + d = 0
\]
Here, $d = -ax_0 - by_0 - cz_0$ and the from above is called
the \textbf{general form}. Continuing, we can establish an
important theorem on the general form.

\begin{theorem}{}{Planes General Form Normal Vector}
For a plane $ax + by + cz + d = 0$ with $\langle a, b, c \rangle
\neq \langle 0, 0, 0 \rangle$, the normal vector is
$\langle a, b, c \rangle$.
\end{theorem}
\begin{proof}
First, consider the case when $a \neq 0$. Rearranging the equation,
\[
ax + by + cz + d
= a \left( x + \frac{d}{a} \right) + by + cz
\]
is obtained and $a \left( x + \frac{d}{a} \right) + by + cz = 0$,
which represents the plane with the normal vector
$\langle a, b, c \rangle$ and passes through the point
$(-\frac{d}{a}, 0, 0)$. The case when $b \neq 0$ or $c \neq 0$
is proven in analogous way. Lastly, if all $a, b, c$ are nonzero,
then it is a plane that is shifted from the origin.
\end{proof}

Let's take a look at quick example.

\begin{problem}{}{}
Find the equation of the plane that contains $\langle 1, 2, 3 \rangle$
and perpendicular to $\langle 1, 2, 3 \rangle$ in point-normal form
and in general form.
\end{problem}
\begin{solution}
Using the form $a(x - x_0) + b(y - y_0) + c(z - z_0) = 0$,
the equation of the plane can be represented as
$(x - 1) + 2(y - 2) + 3(z - 3) = 0$, or $x + 2y + 3z - 14 = 0$.
\end{solution}

Similar to lines, there can be multiple representations of the
same plane. Now before we move on to curves, let's discuss few
meaningful geometric interpretations of planes.

The two cases that we will check are the angle between two planes
and the distance between two parallel planes.

\begin{definition}{}{}
The smaller angle between two planes are known as the \textbf{acute
angle of intersection} between the two planes.
\end{definition}

I know we have the term ``acute'' in it, but $\frac{\pi}{2}$ is
included as acute angle of intersection if the two planes make
a right angle. Now to find an interesting equation like Theorem
\ref{theo:The Dot Product and the Angle}, let's prove a lemma.

\begin{lemma}{}{Intersection Angles Lemma}
For two intersecting lines $L_A$ and $L_B$ in 3-space, let $\theta$
be $\angle(L_A, L_B)$. If $n_A$ and $n_B$ are vectors that are
parallel to the lines respectively, then the following equation holds.
\[
\cos\theta
= \frac{| n_A \cdot n_B |}{\| n_A \| \| n_B \|}
\]
\end{lemma}
\begin{proof}
Let $\alpha$ be the angle between $n_A$ and $n_B$. If $0 < \alpha <
\frac{\pi}{2}$, then the following equation holds by Theorem
\ref{theo:The Dot Product and the Angle}.
\[
\cos\theta
= \cos\alpha
= \frac{n_A \cdot n_B}{\| n_A \| \| n_B \|}
= \frac{| n_A \cdot n_B |}{\| n_A \| \| n_B \|}
\]
Similarly, consider the case when $\frac{\pi}{2} \leq \alpha < \pi$.
By the property of cosine functions, the following equation holds.
\[
\cos\theta
= -\cos\alpha
= |\cos\alpha|
= \frac{| n_A \cdot n_B |}{\| n_A \| \| n_B \|}
\]
Thus, the lemma holds.
\end{proof}

With the lemma shown, we can derive an important theorem regarding
the angle between two intersecting planes.

\begin{theorem}{}{Intersection Angle Theorem}
For two intersecting planes $A$ and $B$ with normal vectors
$n_A$ and $n_B$, if $\theta$ is the acute angle of intersection
between $A$ and $B$, then the following equation holds.
\[
\cos\theta
= \frac{| n_A \cdot n_B |}{\| n_A \| \| n_B \|}
\]
\end{theorem}
\begin{proof}
Let $l_A$ and $l_B$ be two lines normal to planes $A$ and $B$
respectively that also passes through the intersection of the planes.
Let $\alpha$ and $\beta$ be the $\angle(l_A, l_B)$ and
$\angle(l_B, A)$ respectively. Notice that $\alpha + \beta =
\frac{\pi}{2} = \beta + \theta $ through geometric interpretations.
Therefore, $\alpha = \theta$ and it suffices to find $\cos\alpha$.
Because the angle between $l_A$ and $l_B$ and the angle between
$n_A$ and $n_B$ equal, the theorem holds by Lemma
\ref{theo:Intersection Angle Theorem}.
\end{proof}

Continuing, let's take a look at the case when two different planes
do not intersect. To begin, consider the following theorem.

\begin{theorem}{}{Distance Between Plane and Point}
For a plane $ax + by + cz + d = 0$ and a point $(x_0, y_0, z_0)$,
the distance $D$ between the plane and the point is represented with
the following formula.
\[
D
= \frac{| ax_0 + by_0 + cz_0 + d |}{\sqrt{a^2 + b^2 + c^2}}
\]
\end{theorem}
\begin{proof}
Let $Q(x_1, y_1, z_1)$ be any point in the plane. If the angle
between the normal vector $\mathbf{n}$ of the plane from $Q$ and
the vector from $Q$ to $P$ is $\theta$, then the distance $D$ can be
represented as $\| \vec{QP} \| |\cos\theta|$. Thus, the following
equations hold by Theorem \ref{theo:The Dot Product and the Angle}.
\[
D
= \| \vec{QP} \| \cos\theta
= \frac{\| \mathbf{n} \| \| \vec{QP} \| |\cos\theta|}{\| \mathbf{n} \|}
= \frac{| \mathbf{n} \cdot \vec{QP} |}{\| \mathbf{n} \|}
\]
Defining $\mathbf{n} = \langle a, b, c \rangle$, the following
equations hold by construction.
\begin{align*}
    | \mathbf{n} \cdot \vec{QP} |
    &= \left|
        \langle a, b, c \rangle \cdot
            \langle x_0 - x_1, y_0 - y_1, z_0 - z_1 \rangle
    \right| \\
    &= | ax_0 + by_0 + cz_0 + d |
\end{align*}
Because $\| \mathbf{n} \| = \sqrt{a^2 + b^2 + c^2}$, the theorem
satisfy.
\end{proof}

Continuing from the theorem above, we could see that the distance
between two planes is the distance between a plane and a point on
the other plane. Note that both planes $ax + by + cz + d_1 = 0$
and $ax + by + cz + d_2 = 0$ are parallel. Therefore, the distance
$D$ between the two planes can be represented as the following
by Theorem \ref{theo:Distance Between Plane and Point}.
\[
D
= \frac{| ax_0 + by_0 + cz_0 + d_1 |}{\sqrt{a^2 + b^2 + c^2}}
= \frac{| d_1 - d_2 |}{\sqrt{a^2 + b^2 + c^2}}
\]
Notice that this holds because $(x_0, y_0, z_0)$ lies on
$ax + by + cz + d_2 = 0$.

\begin{important}
The distance $D$ between two parallel planes $ax + by + cz + d_1 = 0$
and $ax + by + cz + d_2 = 0$ can be represented as the following.
\[
D
= \frac{| d_1 - d_2 |}{\sqrt{a^2 + b^2 + c^2}}
\]
\end{important}

This is it for this part! In the next part, let's continue from
this linear ideas to expand them to curves.



\end{document}


